\needspace{0.25\textheight}
\section{Отсутствие сингулярностей в модели}\label{sec:noblowup}

\begin{theorem}[главная]\label{th:main}
Пусть выполнена \textup{(H1)}. Тогда для решения \eqref{eq:epd}:
\begin{enumerate}[label=\textup{(\roman*)}, leftmargin=2.4em, itemsep=1pt]
  \item решение существует и единственно при всех $t \ge 0$ при любом
        $\nu \ge 0$ (теорема~\ref{th:wp});
  \item \emph{априорная оценка}: $\displaystyle\sup_{t\ge0}\|\Om(t)\|
        \le \|\Om_0\| < \infty$ (теорема~\ref{th:casimir});
  \item при $\nu = 0$: $\|\Om(t)\| \equiv \|\Om_0\|$, орбита замкнута
        и периодична, если $\Om_0$ не осевая
        (предложение~\ref{prop:periodic});
  \item при $\nu > 0$, $\lambda_{\min}(\bbL) > 0$:
        $\|\Om(t)\| \le \|\Om_0\|\,e^{-\nu\lambda_{\min}t} \to 0$
        (следствие~\ref{cor:decay}, теорема~\ref{th:lasalle});
  \item при \textup{(H2)}: $\mathcal H(t) \le \mathcal H(0)$
        (теорема~\ref{th:ham});
  \item \emph{нет сингулярности}: для любого $T^{*} \le \infty$
        \begin{equation}\label{eq:noblowup}
          \limsup_{t \to T^{*}} \|\Om(t)\| \;\le\; \|\Om_0\| \;<\; \infty ,
        \end{equation}
        и все «серриновские» интегралы конечны тривиально:
        $\int_0^{T}\|\Om(t)\|^{p}dt \le \|\Om_0\|^{p}T < \infty$.
\end{enumerate}
\end{theorem}

\begin{proof}
Пункты (i)--(v)~--- сборка результатов \S\ref{sec:ep}--\S\ref{sec:inv};
покажем (vi). По (ii) функция $t \mapsto \|\Om(t)\|$ ограничена
константой $\|\Om_0\|$ на всём $[0,\infty)$, откуда
$\limsup_{t\to T^{*}}\|\Om(t)\| \le \|\Om_0\|$ для любого
$T^{*} \le \infty$. Конечность интегралов~--- непосредственная
оценка под интегралом.
\end{proof}

\begin{proof}[Комментарий к доказательству]
Содержание оценки \eqref{eq:noblowup}~--- в её механизме: вся
нелинейность $\Om\times\bbI^{-1}\Om$ есть чисто кососимметричный
поворот фазы, ортогональный вектору состояния
($\langle\Om,\Om\times\om\rangle = \det[\Om,\Om,\om] = 0$ при всех
$\om$), а потому \emph{в принципе не способна увеличивать норму}
завихренности; всё возрастание подавлено структурой, а не малостью
коэффициентов. Знак определяет только диссипативный член, и он
неотрицателен по \textup{(H1)}.
\end{proof}

\subsection{Явные решения: осесимметричный случай}

\begin{proposition}[явные решения]\label{prop:axisym}
Пусть $I_1 = I_2 =: I_\perp$ и $\bbL = \diag(\lambda,\lambda,L_{33})$
изотропен в плоскости симметрии. Положим
$z(t) := \Omega_1(t) + i\,\Omega_2(t)$,
$z_0 := \Omega_1(0) + i\,\Omega_2(0)$,
$\Omega_3^0 := \Omega_3(0)$,
$\kappa_0 := \bigl(\tfrac1{I_3}-\tfrac1{I_\perp}\bigr)\Omega_3^0$.
Тогда решение \eqref{eq:epd} выражается в явном виде:
\begin{align}
  \nu = 0:\qquad & \Omega_3(t) \equiv \Omega_3^0, \qquad
    z(t) = z_0\,e^{-i\kappa_0 t}, \label{eq:exact-nu0}\\
  \nu > 0,\ L_{33} > 0:\qquad & \Omega_3(t) = \Omega_3^0\,e^{-\nu L_{33}t},
    \qquad z(t) = z_0\,
    \exp\!\Bigl(-\nu\lambda t - i\,\kappa_0\,
    \frac{1-e^{-\nu L_{33}t}}{\nu L_{33}}\Bigr). \label{eq:exact-nup}
\end{align}
\end{proposition}

\begin{proof}
При $I_1 = I_2$ система \eqref{eq:components} принимает вид
\[
  \dot\Omega_1 = \kappa(t)\,\Omega_2 - \nu\lambda\,\Omega_1,\qquad
  \dot\Omega_2 = -\kappa(t)\,\Omega_1 - \nu\lambda\,\Omega_2,\qquad
  \dot\Omega_3 = -\nu L_{33}\,\Omega_3,
\]
где $\kappa(t) = \bigl(\tfrac1{I_3}-\tfrac1{I_\perp}\bigr)\Omega_3(t)$
(перекрёстные члены $\propto \Omega_1\Omega_3$, $\Omega_2\Omega_3$
симметричны и сокращаются). Третье уравнение~--- линейное, откуда
$\Omega_3(t) = \Omega_3^0e^{-\nu L_{33}t}$ (при $\nu = 0$~---
константа). Для $z$: $\dot z = \dot\Omega_1 + i\dot\Omega_2
= \kappa\Omega_2 - \nu\lambda\Omega_1 + i(-\kappa\Omega_1 - \nu\lambda\Omega_2)
= -i\kappa(t)\,z - \nu\lambda\,z$. Это линейное уравнение с переменным
коэффициентом, и
\[
  z(t) = z_0\exp\!\Bigl(\int_0^t\bigl(-i\kappa(s) - \nu\lambda\bigr)ds\Bigr),
  \qquad
  \int_0^t \kappa(s)\,ds = \kappa_0\,\frac{1-e^{-\nu L_{33}t}}{\nu L_{33}},
\]
что даёт \eqref{eq:exact-nup}; при $\nu = 0$ интеграл равен $\kappa_0 t$,
и получается \eqref{eq:exact-nu0}.
\end{proof}

\begin{remark}[конечная фаза]
Формула \eqref{eq:exact-nup} вскрывает тонкий эффект: суммарный угол
поворота фазы за бесконечное время конечен,
\[
  \int_0^\infty \kappa(s)\,ds = \frac{\kappa_0}{\nu L_{33}} < \infty :
\]
вязкое затухание компоненты $\Omega_3$ останавливает прецессию.
При $\nu = 0$ движение~--- чистая прецессия с периодом
$T = 2\pi/|\kappa_0|$ (частный случай
предложения~\ref{prop:periodic}). Численная проверка точности формул
\eqref{eq:exact-nu0}--\eqref{eq:exact-nup} и порядка сходимости RK4~---
тесты A1--A2 в \S\ref{sec:numerics}.
\end{remark}
